\documentclass[11pt]{article}

\usepackage[margin=1in]{geometry}
\usepackage{amsmath,amssymb,amsfonts}
\usepackage{booktabs}
\usepackage{enumitem}
\usepackage{hyperref}
\usepackage{microtype}
\usepackage[T1]{fontenc}
\usepackage{lmodern}

\hypersetup{
  colorlinks=true,
  linkcolor=blue,
  urlcolor=blue,
  citecolor=blue,
}

\title{Research Journey: From Regret Replay to Learning-Progress Replay}
\author{}
\date{}

\begin{document}
\maketitle

This note records the current mathematical position of the project. The codebase
started as a compact JAX implementation of unsupervised environment design
(UED), with domain randomization, PLR/RPLR, ACCEL, and PAIRED as reference
methods. The current research direction is narrower and more experimental:
we are using the existing PLR machinery as a controlled substrate for asking
which scalar score should cause a level to be replayed.

The central question is:
\begin{quote}
Given a generated level, can we score it by the amount of useful learning it
induces, rather than only by regret-like proxies?
\end{quote}

In the current implementation, the candidate answers are:
\begin{enumerate}[leftmargin=*]
  \item classical PLR regret scores, \texttt{MaxMC} and \texttt{pvl};
  \item absolute raw policy-gradient magnitude, \texttt{abs\_pg};
  \item PPO value-loss magnitude, \texttt{ppo\_value\_loss};
  \item normalized holdout learning progress, \texttt{s\_in}.
\end{enumerate}

The first two are inherited baselines. The latter three represent the present
research frontier.

\section{Shared Setting}

Let $x$ denote a level sampled from the level space $\mathcal{X}$. A policy
$\pi_\theta(a \mid o)$ and value function $V_\theta(o)$ are trained with PPO
on rollouts from levels. For a rollout of horizon $T$ over $N$ parallel
levels, we write:
\begin{itemize}[leftmargin=*]
  \item $o_{t,n}$ for the observation at step $t$ on level $n$;
  \item $a_{t,n}$ for the sampled action;
  \item $r_{t,n}$ for the reward;
  \item $d_{t,n}$ for the episode termination indicator;
  \item $v_{t,n} = V_\theta(o_{t,n})$;
  \item $A_{t,n}$ for the GAE advantage;
  \item $y_{t,n} = A_{t,n} + v_{t,n}$ for the value target.
\end{itemize}

The level sampler maintains a finite buffer
\begin{equation}
  \mathcal{B} = (x_i, s_i, \tau_i)_{i=1}^{m}.
\end{equation}

Here $s_i$ is the current score of level $x_i$, and $\tau_i$ records when the
level was last sampled. Replay weights combine score prioritization and
staleness:
\begin{equation}
  w_i = (1 - \rho)\, w_i^{\mathrm{score}} + \rho\, w_i^{\mathrm{stale}}.
\end{equation}

For rank prioritization, high scores receive low ranks and therefore higher
sampling probability:
\begin{equation}
  w_i^{\mathrm{score}} \propto \operatorname{rank}_i^{-1 / \tau_{\mathrm{temp}}}.
\end{equation}

Thus the research problem reduces to defining the score function
\begin{equation}
  s_\theta(x) \in \mathbb{R}.
\end{equation}

Replay should then emphasize levels that are not merely hard, but useful for
improving the learner.

\section{Baseline Regret Scores}

The existing PLR baselines score a level using regret-like quantities computed
from completed episodes.

\subsection{MaxMC}

Let $R_n^{\max}$ be the maximum Monte Carlo return observed on level $n$. The
MaxMC score is the episode-time average
\begin{equation}
  s_{\mathrm{MaxMC}}(n)
    = \frac{1}{|\mathcal{T}_n|}
      \sum_{t \in \mathcal{T}_n}
      \left(R_n^{\max} - V_\theta(o_{t,n})\right).
\end{equation}

This estimates how far the current value function is below the best return
seen so far. A large score means the level appears to contain unrealized return.

\subsection{Positive Value Loss}

The positive-value-loss score keeps only positive advantages:
\begin{equation}
  s_{\mathrm{PVL}}(n)
    = \frac{1}{|\mathcal{T}_n|}
      \sum_{t \in \mathcal{T}_n}
      \max(A_{t,n}, 0).
\end{equation}

This prioritizes levels where the sampled trajectory performed better than the
critic expected. It is still a regret proxy: it assumes positive surprise is
evidence that the level can teach the agent something.

These scores are sensible baselines, but neither directly asks whether an
update on the level improves generalization, policy fit, or held-out loss.

\section{Candidate 1: Absolute Policy-Gradient Score}

The \texttt{abs\_pg} direction asks whether a level should be replayed when it induces
a large policy update signal. For one timestep and one environment, define the
raw policy-gradient loss
\begin{equation}
  \ell_{\mathrm{pg}}(\theta; t,n)
    = -\log \pi_\theta(a_{t,n} \mid o_{t,n})\, A_{t,n}.
\end{equation}

The per-step score primitive is the Euclidean norm of its parameter gradient:
\begin{equation}
  g_{t,n}
    = \left\lVert \nabla_\theta \ell_{\mathrm{pg}}(\theta; t,n) \right\rVert_2.
\end{equation}

The level score is the time-average of this quantity over completed episodes:
\begin{equation}
  s_{\mathrm{abspg}}(n)
    = \frac{1}{|\mathcal{T}_n|}
      \sum_{t \in \mathcal{T}_n} g_{t,n}.
\end{equation}

Mathematically, this score treats a level as useful when it applies force to
the policy parameters. It is attractive because it is closer to the actual
actor update than MaxMC or PVL. It is also risky because large gradient norm can
mean several different things:
\begin{itemize}[leftmargin=*]
  \item useful policy error;
  \item noisy advantage estimates;
  \item rare high-leverage states;
  \item instability or scale mismatch in the policy parameterization.
\end{itemize}

The implementation therefore normalizes advantages before computing the raw
gradient norms, and computes the score over rollout timesteps rather than from
a single aggregate PPO loss. This gives a matrix $g \in \mathbb{R}^{T \times N}$,
which can be reduced by time to obtain one scalar per level.

The open question is whether gradient magnitude is a useful replay signal by
itself, or whether it needs to be signed, normalized by update variance, or
combined with a held-out progress check.

\section{Candidate 2: PPO Value-Loss Score}

The \texttt{ppo\_value\_loss} direction isolates critic error. For each timestep and
level, define the unclipped value prediction error:
\begin{equation}
  \ell_v(t,n)
    = \frac{1}{2}\left(V_\theta(o_{t,n}) - y_{t,n}\right)^2.
\end{equation}

The level score is
\begin{equation}
  s_{\mathrm{value}}(n)
    = \frac{1}{T}\sum_{t=0}^{T-1}\ell_v(t,n).
\end{equation}

This score prioritizes levels where the critic is poorly calibrated against the
GAE target. It is easier to compute than \texttt{abs\_pg} and is numerically stable,
but it only measures value-function fit. It does not directly say that the
actor improves after replaying the level.

The hypothesis is that critic uncertainty may still be a strong proxy for
useful replay, especially in sparse maze levels where the value function can
identify underlearned structure before the policy-gradient signal is clean.

The main mathematical weakness is that value error can be high for reasons that
do not help the actor: target noise, bootstrapping error, or stochasticity in
episode termination. This score is therefore best viewed as a cheaper baseline
against the more direct learning-progress estimator.

\section{Candidate 3: Normalized Holdout Learning Progress, $S_{\mathrm{in}}$}

The \texttt{s\_in} direction is the closest to the research goal. Instead of asking
whether a level is hard, surprising, or gradient-producing, it asks whether a
small virtual update on one rollout improves PPO loss on an independent rollout
from the same level.

For each level $x_n$, collect two independent rollouts:
\begin{equation}
  D_A(n), D_B(n) \sim \pi_\theta \text{ interacting with } x_n.
\end{equation}

$D_A$ is the update batch. $D_B$ is the holdout evaluation batch. Starting from
the current training state $\theta$, run $K$ virtual PPO updates on $D_A$:
\begin{equation}
  \theta'_n = U_K(\theta, D_A(n)).
\end{equation}

Then evaluate the PPO total loss on $D_B$ before and after the virtual update.
The implementation uses the per-level PPO loss
\begin{equation}
  L_{\mathrm{PPO}}(\theta; D_B)
    = L_{\mathrm{clip}}(\theta; D_B)
      + c_v L_{\mathrm{vf}}(\theta; D_B)
      - c_H H(\theta; D_B).
\end{equation}

This uses the clipped policy objective, clipped value loss, and entropy bonus
matching the normal PPO training objective.

The normalized learning-progress score is
\begin{equation}
  S_{\mathrm{in}}(n)
    = \frac{
      L_{\mathrm{PPO}}(\theta; D_B(n))
      -
      L_{\mathrm{PPO}}(\theta'_n; D_B(n))
    }{
      L_{\mathrm{PPO}}(\theta; D_B(n)) + \epsilon
    }.
\end{equation}

A positive value means that virtual training on $D_A$ reduced held-out loss on
$D_B$. A value near zero means little measurable progress. A negative value
means the virtual update harmed held-out fit.

This score is important because it separates three concepts that earlier
signals conflate:
\begin{itemize}[leftmargin=*]
  \item difficulty: the level currently has high loss or regret;
  \item update magnitude: the level produces a large gradient;
  \item learnability: an update on the level improves independent data from that level.
\end{itemize}

$S_{\mathrm{in}}$ is therefore the most faithful current estimator of ``useful learning.''
Its cost is much higher: each scored level requires independent A/B rollouts,
virtual optimization, and held-out PPO evaluation.

\section{Current Experimental Position}

The project is now positioned to compare replay signals under the same robust
PLR training loop. The relevant score functions are:
\begin{center}
\begin{tabular}{@{}ll@{}}
\toprule
Score function & Description \\
\midrule
\texttt{MaxMC} & baseline regret against best observed return \\
\texttt{pvl} & baseline positive advantage / positive value loss \\
\texttt{abs\_pg} & raw actor-gradient magnitude \\
\texttt{ppo\_value\_loss} & critic target error \\
\texttt{s\_in} & held-out PPO loss reduction after virtual updates \\
\bottomrule
\end{tabular}
\end{center}

The default experimental setting in the current run scripts is robust PLR on
maze levels with:
\begin{center}
\begin{tabular}{@{}ll@{}}
\toprule
Parameter & Value \\
\midrule
\texttt{num\_train\_envs} & 32 \\
\texttt{num\_steps} & 256 \\
\texttt{num\_updates} & 30000 \\
\texttt{replay\_prob} & 0.8 \\
\texttt{level\_buffer\_capacity} & 4000 \\
\texttt{prioritization} & rank \\
\texttt{staleness\_coeff} & 0.3 \\
\texttt{n\_walls} & 60 \\
\bottomrule
\end{tabular}
\end{center}

Evaluation is against the fixed maze suite:
\begin{center}
\texttt{SixteenRooms}, \texttt{SixteenRooms2}, \texttt{Labyrinth}, \texttt{LabyrinthFlipped},\\
\texttt{Labyrinth2}, \texttt{StandardMaze}, \texttt{StandardMaze2}, \texttt{StandardMaze3}.
\end{center}

So the journey has moved from ``implement standard UED algorithms'' to ``use PLR
as an experimental harness for learning-progress-aware level selection.''

\section{What Is Already Mathematically Settled}

The following pieces are now explicit enough to treat as fixed for the next
round of experiments:
\begin{enumerate}[leftmargin=*]
  \item The PLR buffer mechanics are not the research variable. They provide the replay substrate.
  \item \texttt{MaxMC} and \texttt{pvl} are the regret-style controls.
  \item \texttt{abs\_pg} defines actor-side update magnitude at the per-timestep level.
  \item \texttt{ppo\_value\_loss} defines critic-side prediction error as a per-level score.
  \item \texttt{s\_in} defines a direct holdout learning-progress estimator.
  \item All candidate scores reduce to one scalar per level, allowing the same sampler to rank, insert, replay, and update levels.
\end{enumerate}

This is a useful point in the journey because the score definitions now share a
common interface while expressing different theories of what replay should
mean.

\section{What Remains Open}

The research questions are now empirical and interpretive rather than merely
implementational.

\subsection{Does high gradient norm predict improvement?}

\texttt{abs\_pg} may identify levels with strong policy signal, but it may also
prioritize noise. A useful diagnostic is the correlation between
$s_{\mathrm{abspg}}(n)$
and subsequent improvement in evaluation return or $S_{\mathrm{in}}(n)$.

\subsection{Is critic error enough?}

\texttt{ppo\_value\_loss} is cheap and stable. The question is whether critic error is
aligned with useful policy learning. If it performs close to $S_{\mathrm{in}}$, then the
expensive virtual-update estimator may not be necessary.

\subsection{Does $S_{\mathrm{in}}$ justify its compute cost?}

$S_{\mathrm{in}}$ is conceptually strongest but computationally expensive. Its value
depends on whether it produces better level curricula than cheaper proxies. The
key comparison is not only final return, but sample efficiency and replay
composition over training.

\subsection{Should negative learning progress be retained?}

The current formula permits negative scores:
\begin{equation}
  S_{\mathrm{in}}(n) < 0
\end{equation}
when a virtual update worsens held-out PPO loss. This is mathematically
meaningful, but the replay sampler ranks scores. We need to decide whether
negative progress should suppress replay, be clipped to zero, or be treated as
a warning signal for overly hard or unstable levels.

\subsection{Should scores be normalized across training time?}

The scale of gradients, value errors, and PPO losses changes throughout
training. Rank prioritization reduces some scale sensitivity, but cross-time
comparisons still matter for buffer insertion and replacement. A natural next
step is to track score distributions over time and consider running
normalization or quantile-based insertion rules.

\section{Near-Term Research Plan}

The next mathematical milestone is to turn the score definitions into a clean
comparison table:
\begin{center}
\begin{tabular}{@{}p{0.22\linewidth}p{0.34\linewidth}p{0.34\linewidth}@{}}
\toprule
Score function & What it estimates & Expected failure mode \\
\midrule
\texttt{MaxMC} & unrealized return & stale best-return memory \\
\texttt{pvl} & positive surprise & advantage noise \\
\texttt{abs\_pg} & actor update magnitude & noisy/high-variance gradients \\
\texttt{ppo\_value\_loss} & critic prediction error & value-only misalignment \\
\texttt{s\_in} & held-out learning progress & high compute cost \\
\bottomrule
\end{tabular}
\end{center}

The next experimental milestone is to run controlled robust-PLR sweeps where
only the score function changes. The primary comparisons should be:
\begin{enumerate}[leftmargin=*]
  \item final evaluation return on the fixed maze suite;
  \item learning speed as a function of environment steps;
  \item buffer composition over time;
  \item score distribution over time;
  \item correlation between cheap proxies and $S_{\mathrm{in}}$.
\end{enumerate}

If \texttt{abs\_pg} or \texttt{ppo\_value\_loss} correlates strongly with $S_{\mathrm{in}}$, then the
project can move toward a cheaper approximation of learning-progress replay.
If not, then $S_{\mathrm{in}}$ becomes the principled target, and the next problem is to
make it cheaper through subsampling, fewer virtual updates, cached rollouts, or
less frequent scoring.

\section{Current Thesis}

The working thesis is:
\begin{quote}
Regret-based PLR replays levels that appear hard, but the better replay rule
is to prioritize levels whose data induces measurable learning progress under
the current optimizer.
\end{quote}

The present codebase has reached the point where this thesis is testable. The
mathematical objects are defined, the score functions share the same replay
interface, and the remaining work is to determine which score is predictive
enough to guide curriculum formation in practice.

\end{document}
